\section{Introduction}
\label{sec:intro}

Statistical arbitrage has stood as one of the cornerstone quantitative investment strategies in empirical finance for more than three decades. Originating in Wall Street proprietary trading desks during the mid-1980s, the core economic hypothesis posits that assets sharing common fundamental, statistical, or factor-driven risks should display cointegrated or mean-reverting price trajectories \citep{gatev2006pairs,vidyamurthy2004pairs}. When temporary idiosyncratic supply-demand imbalances or liquidity shocks drive prices apart, statistical arbitrageurs inject liquidity into divergent assets by buying undervalued securities and shorting overvalued counterparts, expecting price convergence to generate risk-adjusted returns \citep{huck2015pairs,avellaneda2010statistical,krauss2017statistical,gu2020empirical,guijarro2021deep}.

However, modern empirical finance has increasingly demonstrated that anomalies documented in frictionless backtests often deteriorate drastically out-of-sample \citep{mclean2016does,hou2020replicating}. In the specific domain of pairs trading, \citet{do2010does} and \citet{do2012are} discovered that post-2002 US pairs trading profits declined steadily, with trading friction and short-selling borrowing constraints consuming the majority of remaining gains. As \citet{novymarx2016taxonomy} and \citet{frazzini2018trading} have shown across broader asset pricing anomalies, strategies characterized by high portfolio turnover are acutely susceptible to transaction cost erosion.

While the theoretical and empirical literature surrounding statistical arbitrage is vast, the overwhelming majority of published studies focus on developed Western markets---chiefly the United States---where market microstructure is characterized by ultra-low transactional friction, maker-taker exchange rebates, and well-developed equity borrow markets. Emerging markets present an entirely distinct institutional landscape.

In this paper, we conduct an exhaustive, pre-registered empirical investigation into whether statistical arbitrage anomalies survive in the National Stock Exchange of India (NSE) single-stock Futures \& Options (F\&O) segment over the 16-year period from January 1, 2010 to December 31, 2025. The Indian equity derivatives market serves as a critical empirical laboratory for two institutional reasons:
\begin{enumerate}
    \item It represents one of the largest and most active derivative venues globally by contract volume, yet single-stock short selling in the cash equity segment is severely constrained by illiquid Securities Lending and Borrowing (SLB) mechanisms, forcing market-neutral managers to operate almost exclusively in single-stock futures contracts.
    \item Unlike US equities, Indian derivative transactions are subject to non-negotiable statutory government taxes, including the Securities Transaction Tax (STT) \citep{bose2007securities}, Stamp Duty, Goods and Services Tax (GST), exchange turnover levies, and SEBI regulatory charges.
\end{enumerate}

\subsection{Key Empirical Findings and the ``Gross Alpha Illusion''}
Our central finding is striking: \textbf{Gross statistical arbitrage signals exist in Indian F\&O equities with high statistical significance, but they are completely annihilated by market microstructure friction.}

Across all three methodological families examined:
\begin{itemize}
    \item \textbf{Strategy A (Gatev et al.\ 2006 Distance Pairs):} Generates a gross annualized Sharpe ratio of $+0.95$ ($t\text{-stat} = +4.11$, $p < 0.001$). Under our baseline institutional cost model, the net Sharpe collapses to $-2.52$ with an annualized net loss of $-25.16\%$.
    \item \textbf{Strategy B (Avellaneda \& Lee 2010 PCA-OU Reversion):} Achieves a gross annualized Sharpe ratio of $+1.74$ ($t\text{-stat} = +7.49$, $p < 0.001$). After dated transaction costs, the net Sharpe collapses to $-1.63$ ($t\text{-stat} = -7.00$), yielding a $-16.75\%$ net loss per annum.
    \item \textbf{Strategy C (ML-Enhanced Walk-Forward Residuals):} Demonstrates an impressive gross Sharpe of $+1.00$ ($t\text{-stat} = +4.05$, $p < 0.001$). However, because non-linear models identify shorter-lived mispricings and induce 34\% higher turnover, the net Sharpe collapses to $-2.96$ ($t\text{-stat} = -12.00$, net loss $-25.91\%$).
\end{itemize}

The breakeven cost multiplier across all strategies averages $0.38\times$ of baseline institutional costs. Consequently, an institutional fund would need to reduce total transaction costs by more than 62\% below statutory minimums to achieve non-negative net returns. Controlling for data-snooping via the Deflated Sharpe Ratio \citep{bailey2014deflated} and multiple testing \citep{harvey2016and,benjamini1995controlling}, zero out of twelve pre-registered strategy specifications survive. We conclude that daily statistical arbitrage in Indian single-stock F\&O is not viable as an independent, standalone alpha generator, but retains substantial economic value when repurposed as an execution timing overlay \citep{novymarx2016taxonomy}.
